\clearpage
\appendix
\numberwithin{figure}{section}
\numberwithin{table}{section}
\setcounter{figure}{0}
\setcounter{table}{0}
\renewcommand{\thefigure}{\Alph{section}.\arabic{figure}}
\section{Distribution and source-gap evidence}\label{app:data}
These supplementary figures separate the distribution of measured rain from the absence of measurements. They support the data-quality discussion without treating an unrecorded interval as dry weather.
\plot{02_positive_distribution}{Distribution of positive five-minute rainfall across 2017--2024. Bars give the share of positive observations in log-spaced bins; the horizontal axis is logarithmic. Recorded zeros are excluded. The long upper tail motivates evaluating rare accumulations separately.}{fig:appdistribution}
\plot{04_source_gap}{Daily network-wide raw and cleaned record counts in August 2018. The absence of observations on 4--31 August is already present in the source file. The cleaning process therefore does not explain this gap, and this figure does not identify its upstream cause.}{fig:appgap}
\FloatBarrier

\clearpage
\section{Monthly trajectories and annual support}\label{app:temporal}
The monthly trajectories retain usable periods individually. The annual plot applies the stricter full-year rule, showing why monthly evidence is more extensive than annual evidence.
\plot{07_monthly_trends}{Mean observed monthly station totals, with a separate trajectory for each calendar month. Missing year-months break the lines. Vertical ranges differ across panels; use the main monthly heatmap to compare absolute amounts between months. Station composition can change between points, so these are descriptive network trajectories.}{fig:apptrajectories}
\plot{08_annual_totals}{Observed calendar-year totals at eligible stations. Points show station totals, coral diamonds their mean, and $n$ the contributing station count. Grey columns have no station meeting annual coverage of at least 95\% and coverage of at least 90\% in every month. The three supported years use different station sets and do not establish a national wettest-year ranking.}{fig:appannual}
\FloatBarrier

\clearpage
\section{Daily rainfall and coverage sensitivity}\label{app:robustness}
Daily upper-tail amounts and wet-day frequencies describe different aspects of rainfall. The sensitivity curves show how changing the completeness rule affects seasonal summaries, while retaining the limitations of the observed network.
\plot{11_daily_tails}{Monthly daily-rainfall P99 (upper) and wet station-day percentage (lower). Included days have at least 95\% daily coverage within eligible station-months. Quantiles include dry days; wet means observed daily rainfall above zero. P99 is an approximate Spark quantile of station-days, not an independent-storm return level. Grey cells lack qualifying support.}{fig:apptails}
\plot{12_sensitivity}{Seasonal mean amount and wet-interval frequency under 80\%, 90\% and 95\% monthly coverage cutoffs. Each curve uses its own eligible stations and equally weighted available years. Similar curves support the stability of these summaries under the tested cutoffs; they do not establish random missingness.}{fig:appsensitivity}
\FloatBarrier

\clearpage
\section{Additional spatial views}\label{app:spatial}
The twelve-month view complements the main December comparison. Regional colours summarise contributing gauges and should be read together with their locations and coverage.
\plot{14_spatial_year_2024}{All twelve months of 2024. Colours show observed monthly totals at qualified stations, and $n$ gives their number. The original gallery's global monthly colour scale is retained; it differs from the December-only scale in the main text. The eligible network changes by month, and no rainfall is interpolated between gauges.}{fig:appspatialseason}
\plot{15_regions_month_12}{December planning-area station means across eight years. Each coloured polygon shows the arithmetic mean of its assigned eligible stations; dots show contributing locations. Grey means no supporting station, not zero rainfall. The original regional scale is retained. These are gauge summaries, not area-weighted rainfall estimates; station counts and placements affect their representativeness.}{fig:appregions}
\FloatBarrier

\clearpage
\section{Additional predictive diagnostics}\label{app:prediction}
These figures describe the 15\% RF and logistic models on the 2024 test set. They show how performance varies by rainfall threshold and month, alongside the RF's feature use.
\plot{19_baseline_skill}{Brier skill relative to each mode's historical-frequency baseline on its own eligible 2024 test population. Positive values indicate lower error than the baseline; zero indicates equal error. A's random forest improves on the baseline at every evaluated threshold, whereas A's logistic regression deteriorates at larger thresholds. B's forest has small, threshold-dependent gains and losses. The two panels use different eligible populations.}{fig:appbaselines}
\plot{21_diagnostics}{Left: A's monthly Brier skill relative to historical frequency at 1 mm in 2024. Right: its six largest random-forest impurity importances. Recent rainfall totals and persistence summaries contribute most to the fitted forest; these correlated features share information.}{fig:appdiagnostics}
\FloatBarrier

\clearpage
\section{Model scores and experiment coverage}\label{app:experiments}
Table~\ref{tab:allmodels} brings the recorded model and training-fraction results together. Mean Brier and mean log loss give equal weight to the seven primary thresholds; PR-AUC is shown at 1 mm. Historical frequency uses all eligible training examples in every run.
\begin{table}[htbp]\centering\small
\begin{tabular}{llrrrr}\toprule
Mode & Model & Training & Mean Brier & Mean log loss & PR-AUC at 1 mm\\\midrule
@@EXPERIMENT_METRIC_ROWS@@
\bottomrule\end{tabular}
\caption{2024 results on each mode's eligible population: 5,729,508 anchors for A and 6,374,476 for B. The 15\% A GBT row uses its recorded aggregate summary; the 15\% B GBT row uses per-threshold exports. All other rows use saved metric CSVs.}\label{tab:allmodels}
\end{table}

The RF/LR 15\% and 100\% runs retain both model pipelines and test metrics. GBT 30\% and 100\% retain pipelines, daily errors and reliability tables; the 15\% results are retained as summaries. RF's 30\% experiment is evaluated on 2022 validation in Appendix~\ref{sec:rfsensitivity}. The archived logistic runs use 15\% and 100\% training. Within the GBT exports, RF and logistic columns refer to the loaded 15\% comparison models, including at the additional 3, 7.5 and 15 mm queries. The repository's Part 2 guide maps these results to their source directories and query commands.
\FloatBarrier
